% BHU PYQ -- Manually transcribed & error-corrected
% Paper: GRA/GRB-506: Population Geography (Elective)
% Exam: B.A./B.Sc. (Hons.) Semester V Examination, 2022-23

\documentclass[11pt,a4paper]{article}
\usepackage[a4paper,top=2.5cm,bottom=2.5cm,left=2.8cm,right=2.8cm,headheight=14pt]{geometry}
\usepackage{amsmath,amssymb}
\usepackage[T1]{fontenc}
\usepackage[utf8]{inputenc}
\usepackage{lmodern,microtype,fancyhdr,enumitem,hyperref,lastpage,parskip}

\hypersetup{colorlinks=true,urlcolor=black,linkcolor=black,
  pdftitle={GRA/GRB-506 Population Geography -- BHU B.A./B.Sc. Sem V 2022-23}}
\pagestyle{fancy}\fancyhf{}
\renewcommand{\headrulewidth}{0.4pt}\renewcommand{\footrulewidth}{0.4pt}
\fancyhead[L]{\small   \textit{Banaras Hindu University}}
\fancyhead[R]{\small   \textit{GRA/GRB-506: Population Geography}}
\fancyfoot[C]{\small Page  \thepage\ of \pageref{LastPage}}


\newlist{parts}{enumerate}{2}
\setlist[parts,1]{label=(\alph*),leftmargin=*,itemsep=5pt,topsep=3pt}

\newcommand{\pts}[1]{\hfill   \textit{[#1]}}


\begin{document}

\begin{center}
  {\large\bfseries BANARAS HINDU UNIVERSITY}\\[2pt]
  {
\normalsize B.A./B.Sc. (Hons.) --- Semester V Examination, 2022-23}\\[6pt]
  \rule{0.8\linewidth}{0.5pt}\\[6pt]
  {\large\bfseries Paper No. GRA/GRB-506: Population Geography (Elective)}\\[4pt]
  {
\normalsize Time: 3 Hours \hfill Maximum Marks: 70}\\[2pt]
  {\small Ref. No.: 052714}
\end{center}
\rule{\linewidth}{0.4pt}
\smallskip



\noindent   \textit{Instructions: Answer   \textbf{five} questions in all. Question No. 1, carrying 20 marks, is compulsory. Remaining questions are of equal marks.}
\bigskip




\noindent   \textbf{Question 1.} Write short answer of the following: \pts{20}

\begin{parts}
  \item Vital Registration \pts{4}
  \item Migration types \pts{4}
  \item Demographic transition theory \pts{4}
  \item Population problems and planning \pts{4}
  \item Significance of age and sex structure \pts{4}
\end{parts}

\medskip


\noindent   \textbf{Question 2.} Explain the concepts of systematic and regional approaches in studying population geography. \pts{12.5}

\medskip


\noindent   \textbf{Question 3.} Discuss the concept of urbanization. \pts{12.5}

\medskip


\noindent   \textbf{Question 4.} Explain the causes and consequences of world population growth. \pts{12.5}

\medskip


\noindent   \textbf{Question 5.} Describe the major sources of population data. \pts{12.5}

\medskip


\noindent   \textbf{Question 6.} Explain the determinants of fertility. \pts{12.5}

\medskip


\noindent   \textbf{Question 7.} Discuss the indicators and patterns of human resource development in the world. \pts{12.5}

\medskip


\noindent   \textbf{Question 8.} Define population density and discuss the relative significance of various density types. \pts{12.5}

\medskip


\noindent   \textbf{Question 9.} Give an account of the 'Population Policy 2000' of India. \pts{12.5}

\vfill\rule{\linewidth}{0.4pt}

\begin{center}\small
\href{https://bhu-exam.vercel.app/}{Click here to visit our website}\\[4pt]
\href{https://me-aryan.vercel.app/}{Click here to visit our contributor}\\[4pt]
\href{https://quashan-me.vercel.app/}{Click here to visit our contributor}
\end{center}
\end{document}
